% ============================================
% Supplementary Figures
% ============================================
\newpage
\subsection*{Supplementary Figures}

% ============================================
% Section 3.1: CLS is prevalent - Strain Library
% ============================================

% Supple Fig S1 - Taxonomy color map
\begin{figure}[H]
\centering
\includegraphics[width=0.9\textwidth]{supplementary_figs/taxonomy_color_map.pdf}
\caption{Table of phylogenetic taxonomy for 50 ASV isolates used in the experiment. Each row shows the ASV identifier and its full taxonomic classification (Kingdom, Phylum, Class, Order, Family, Genus). Colors match those used in pie charts and composition figures throughout the manuscript (inferno colormap with seed=4).}
\label{fig:suppl_s1}
\end{figure}

% Supple Fig S2 - Phylogeny tree
\begin{figure}[H]
\centering
\includegraphics[width=0.9\textwidth]{supplementary_figs/Fig_1-3_phylogeny_tree.pdf}
\caption{Phylogeny map (EMBL).}
\label{fig:suppl_s2}
\end{figure}

% Supple Fig S3 - Time series in Base media
\begin{figure}[H]
\centering
\begin{subfigure}[b]{0.32\textwidth}
    \includegraphics[width=\textwidth]{supplementary_figs/timeseries_merged/Fig_1-3_timeseries_Base_ex1_merged.pdf}
    \caption{Example 1}
\end{subfigure}
\hfill
\begin{subfigure}[b]{0.32\textwidth}
    \includegraphics[width=\textwidth]{supplementary_figs/timeseries_merged/Fig_1-3_timeseries_Base_ex2_merged.pdf}
    \caption{Example 2}
\end{subfigure}
\hfill
\begin{subfigure}[b]{0.32\textwidth}
    \includegraphics[width=\textwidth]{supplementary_figs/timeseries_merged/Fig_1-3_timeseries_Base_ex3_merged.pdf}
    \caption{Example 3}
\end{subfigure}
\caption{Time series of community composition during coalescence in Base medium. Each panel shows two parental subcommunities (left, stacked vertically) and their coalesced outcome (right) over serial dilution cycles. Stacked bar charts represent relative species abundances at each time point.}
\label{fig:suppl_s3}
\end{figure}

% Supple Fig S4 - Robustness of Fig.1
\begin{figure}[H]
\centering
% File: supplementary_figs/Fig_robustness_comparison_MN_merged.svg (convert to PDF)
\fbox{\parbox{0.9\textwidth}{\centering\vspace{2cm}\textbf{[SVG] Fig\_robustness\_comparison\_MN\_merged.svg}\vspace{2cm}}}
\caption{Robustness of Fig.~1 across various similarity metrics.}
\label{fig:suppl_s4}
\end{figure}

% ============================================
% Section 3.1: Skewness Null Model Analysis
% ============================================

% Supple Fig S5 - Rank-Abundance Curves - Base medium (parental vs coalesced)
\begin{figure}[H]
\centering
\includegraphics[width=0.85\textwidth]{supplementary_figs/rank_abundance_parental_vs_coalesced_Base.pdf}
\caption{Rank-abundance curves comparing parental and coalesced communities in Base medium. (A) Parental (sub)communities before coalescence. (B) Coalesced communities after coalescence. Each line represents one community's rank-abundance curve. Gini coefficients quantify abundance inequality (0 = perfectly even, 1 = one species dominates). Parental: Gini = 0.62 $\pm$ 0.19 (n=59); Coalesced: Gini = 0.64 $\pm$ 0.15 (n=94). Both parental and coalesced communities exhibit substantial abundance skewness with similar Gini values.}
\label{fig:suppl_s5}
\end{figure}

% Supple Fig S6 - Rank-Abundance Curves - Nutr- medium (parental vs coalesced)
\begin{figure}[H]
\centering
\includegraphics[width=0.85\textwidth]{supplementary_figs/rank_abundance_parental_vs_coalesced_Nutr-.pdf}
\caption{Rank-abundance curves comparing parental and coalesced communities in Nutr$-$ medium. (A) Parental (sub)communities before coalescence. (B) Coalesced communities after coalescence. Each line represents one community's rank-abundance curve. Parental: Gini = 0.63 $\pm$ 0.09 (n=60); Coalesced: Gini = 0.61 $\pm$ 0.12 (n=92).}
\label{fig:suppl_s6}
\end{figure}

% Supple Fig S7 - Rank-Abundance Curves - Nutr+ medium (parental vs coalesced)
\begin{figure}[H]
\centering
\includegraphics[width=0.85\textwidth]{supplementary_figs/rank_abundance_parental_vs_coalesced_Nutr+.pdf}
\caption{Rank-abundance curves comparing parental and coalesced communities in Nutr+ medium. (A) Parental (sub)communities before coalescence. (B) Coalesced communities after coalescence. Each line represents one community's rank-abundance curve. Parental: Gini = 0.64 $\pm$ 0.17 (n=60); Coalesced: Gini = 0.66 $\pm$ 0.14 (n=94).}
\label{fig:suppl_s7}
\end{figure}

% Supple Fig S8 - Skewness Null Model Comparison
\begin{figure}[H]
\centering
\includegraphics[width=0.45\textwidth]{supplementary_figs/skewness_null_comparison.pdf}
\caption{Testing whether skewed parental abundance distributions explain one-sided community-level selection. Scatter plots show 100 randomly sampled data points with mean $\pm$ SEM (squares with error bars) for experimental one-sided selection values compared against two null models: (1) abundance-weighted random selection, where species survival probability is proportional to their abundance in the combined parental pool, and (2) shuffled abundance, where abundances are randomly permuted among species within each parental community before neutral mixing. Both null models produce significantly lower one-sided selection than experimental data (Mann-Whitney U test, $p < 0.001$ for both comparisons), indicating that skewed abundance distributions alone do not fully account for the observed asymmetric outcomes.}
\label{fig:suppl_s8}
\end{figure}

% ============================================
% Section 3.3: Simulation Robustness
% ============================================

% Supple Fig S9 - Simulation: growth-rate heterogeneity
\begin{figure}[H]
\centering
\fbox{\parbox{0.9\textwidth}{\centering\vspace{2cm}\textbf{[MISSING] Simulation: growth-rate heterogeneity}\vspace{2cm}}}
\caption{Simulation: growth-rate heterogeneity.}
\label{fig:suppl_s9}
\end{figure}

% Supple Fig S10 - Simulation: carrying-capacity variation
\begin{figure}[H]
\centering
\fbox{\parbox{0.9\textwidth}{\centering\vspace{2cm}\textbf{[MISSING] Simulation: carrying-capacity variation}\vspace{2cm}}}
\caption{Simulation: carrying-capacity variation.}
\label{fig:suppl_s10}
\end{figure}

% Supple Fig S11 - Simulation: alternative alpha_ij distributions
\begin{figure}[H]
\centering
\begin{subfigure}[b]{0.48\textwidth}
    \includegraphics[width=\textwidth]{supplementary_figs/Fig_phase_diagram_ablation_gaussian.pdf}
    \caption{Gaussian distribution: $\alpha_{ij} \sim \mathcal{N}(\mu, (\mu/\sqrt{3})^2)$}
\end{subfigure}
\hfill
\begin{subfigure}[b]{0.48\textwidth}
    \includegraphics[width=\textwidth]{supplementary_figs/Fig_phase_diagram_ablation_gamma.pdf}
    \caption{Gamma distribution: $\alpha_{ij} \sim \text{Gamma}(k=3, \theta=\mu/3)$}
\end{subfigure}
\caption{Phase diagrams for alternative $\alpha_{ij}$ sampling distributions. Both distributions have the same coefficient of variation (CV $= 1/\sqrt{3} \approx 0.577$) but different shapes. (A) Gaussian distribution allows negative interaction coefficients at low mean values. (B) Gamma distribution ensures all coefficients are positive. The qualitative transition from mixing-dominated to dominance/restructuring outcomes with increasing mean interaction strength $\mu$ is robust to the choice of distribution, confirming that the observed patterns depend primarily on interaction strength rather than distributional details.}
\label{fig:suppl_s11}
\end{figure}

% ============================================
% Section 3.4: Pairwise Invasion Assays
% ============================================

% Supple Fig S12 - Nutr- pairwise invasion
\begin{figure}[H]
\centering
\includegraphics[width=0.9\textwidth]{supplementary_figs/Pairwise_Matrix_Dynamics_LN_improved.pdf}
\caption{Nutr$-$ pairwise invasion.}
\label{fig:suppl_s12}
\end{figure}

% Supple Fig S13 - Base pairwise invasion
\begin{figure}[H]
\centering
\includegraphics[width=0.9\textwidth]{supplementary_figs/Pairwise_Matrix_Dynamics_MN_improved.pdf}
\caption{Base pairwise invasion.}
\label{fig:suppl_s13}
\end{figure}

% Supple Fig S14 - Nutr+ pairwise invasion
\begin{figure}[H]
\centering
\includegraphics[width=0.9\textwidth]{supplementary_figs/Pairwise_Matrix_Dynamics_HN_improved.pdf}
\caption{Nutr+ pairwise invasion.}
\label{fig:suppl_s14}
\end{figure}

% ============================================
% Coalescence Outcome Matrices (Supplementary only - not referenced in main)
% ============================================

% Supple Fig S15 - Coalescence Matrix - Nutr- (LN)
\begin{figure}[H]
\centering
\begin{subfigure}[b]{0.24\textwidth}
    \includegraphics[width=\textwidth]{supplementary_figs/PieCharts/CoalescenceMatrices/coalescence_matrix_Nutr-_s6.png}
    \caption{Pool size 6}
\end{subfigure}
\hfill
\begin{subfigure}[b]{0.24\textwidth}
    \includegraphics[width=\textwidth]{supplementary_figs/PieCharts/CoalescenceMatrices/coalescence_matrix_Nutr-_s12.png}
    \caption{Pool size 12}
\end{subfigure}
\hfill
\begin{subfigure}[b]{0.24\textwidth}
    \includegraphics[width=\textwidth]{supplementary_figs/PieCharts/CoalescenceMatrices/coalescence_matrix_Nutr-_s24.png}
    \caption{Pool size 24}
\end{subfigure}
\caption{Coalescence outcome matrices in Nutr$-$ medium. Each matrix shows pairwise coalescence outcomes for all combinations of parental communities. Upper triangle cells display pie charts of the coalesced community composition.}
\label{fig:suppl_s15}
\end{figure}

% Supple Fig S16 - Coalescence Matrix - Base (MN)
\begin{figure}[H]
\centering
\begin{subfigure}[b]{0.24\textwidth}
    \includegraphics[width=\textwidth]{supplementary_figs/PieCharts/CoalescenceMatrices/coalescence_matrix_Base_s6.png}
    \caption{Pool size 6}
\end{subfigure}
\hfill
\begin{subfigure}[b]{0.24\textwidth}
    \includegraphics[width=\textwidth]{supplementary_figs/PieCharts/CoalescenceMatrices/coalescence_matrix_Base_s12.png}
    \caption{Pool size 12}
\end{subfigure}
\hfill
\begin{subfigure}[b]{0.24\textwidth}
    \includegraphics[width=\textwidth]{supplementary_figs/PieCharts/CoalescenceMatrices/coalescence_matrix_Base_s24.png}
    \caption{Pool size 24}
\end{subfigure}
\caption{Coalescence outcome matrices in Base medium. Each matrix shows pairwise coalescence outcomes for all combinations of parental communities. Upper triangle cells display pie charts of the coalesced community composition.}
\label{fig:suppl_s16}
\end{figure}

% Supple Fig S17 - Coalescence Matrix - Nutr+ (HN)
\begin{figure}[H]
\centering
\begin{subfigure}[b]{0.24\textwidth}
    \includegraphics[width=\textwidth]{supplementary_figs/PieCharts/CoalescenceMatrices/coalescence_matrix_Nutr+_s6.png}
    \caption{Pool size 6}
\end{subfigure}
\hfill
\begin{subfigure}[b]{0.24\textwidth}
    \includegraphics[width=\textwidth]{supplementary_figs/PieCharts/CoalescenceMatrices/coalescence_matrix_Nutr+_s12.png}
    \caption{Pool size 12}
\end{subfigure}
\hfill
\begin{subfigure}[b]{0.24\textwidth}
    \includegraphics[width=\textwidth]{supplementary_figs/PieCharts/CoalescenceMatrices/coalescence_matrix_Nutr+_s24.png}
    \caption{Pool size 24}
\end{subfigure}
\caption{Coalescence outcome matrices in Nutr+ medium. Each matrix shows pairwise coalescence outcomes for all combinations of parental communities. Upper triangle cells display pie charts of the coalesced community composition.}
\label{fig:suppl_s17}
\end{figure}

% ============================================
% Additional Time Series (Supplementary only - not referenced in main)
% ============================================

% Supple Fig S18 - Nutr- time series
\begin{figure}[H]
\centering
\begin{subfigure}[b]{0.45\textwidth}
    \includegraphics[width=\textwidth]{supplementary_figs/timeseries_merged/Fig_1-3_timeseries_Nutr-_ex1_merged.pdf}
    \caption{Example 1}
\end{subfigure}
\hfill
\begin{subfigure}[b]{0.45\textwidth}
    \includegraphics[width=\textwidth]{supplementary_figs/timeseries_merged/Fig_1-3_timeseries_Nutr-_ex2_merged.pdf}
    \caption{Example 2}
\end{subfigure}
\caption{Time series of community composition during coalescence in Nutr$-$ medium. Each panel shows two parental subcommunities (left, stacked vertically) and their coalesced outcome (right) over serial dilution cycles. Stacked bar charts represent relative species abundances at each time point.}
\label{fig:suppl_s18}
\end{figure}

% Supple Fig S19 - Nutr+ time series
\begin{figure}[H]
\centering
\includegraphics[width=0.5\textwidth]{supplementary_figs/timeseries_merged/Fig_1-3_timeseries_Nutr+_ex1_merged.pdf}
\caption{Time series of community composition during coalescence in Nutr+ medium. The panel shows two parental subcommunities (left, stacked vertically) and their coalesced outcome (right) over serial dilution cycles. Stacked bar charts represent relative species abundances at each time point.}
\label{fig:suppl_s19}
\end{figure}

% ============================================
% Unreferenced Figures (Need placement suggestions)
% ============================================

% Supple Fig S20 - Parental Dominance Index null model (Base)
\begin{figure}[H]
\centering
% File: supplementary_figs/correlation_barplot_MN.svg (convert to PDF)
\fbox{\parbox{0.9\textwidth}{\centering\vspace{2cm}\textbf{[SVG] correlation\_barplot\_MN.svg}\vspace{2cm}}}
\caption{Parental Dominance Index (PDI) of experiment and null model in Base media.}
\label{fig:suppl_s20}
\end{figure}

% Supple Fig S21 - Pairwise selection correlation (Nutr-/Nutr+)
\begin{figure}[H]
\centering
\begin{subfigure}[b]{0.48\textwidth}
    \includegraphics[width=\textwidth]{supplementary_figs/correlation_barplot_LN.png}
    \caption{Nutr$-$ medium}
\end{subfigure}
\hfill
\begin{subfigure}[b]{0.48\textwidth}
    \includegraphics[width=\textwidth]{supplementary_figs/correlation_barplot_HN.png}
    \caption{Nutr+ medium}
\end{subfigure}
\caption{\textbf{Pairwise selection correlation in experimental coalescence across nutrient conditions.} Mean pairwise selection correlation for species pairs from the same parental community (red) versus cross-parent pairs (blue). Grey horizontal line indicates random selection baseline. Individual dots show per-event correlations; squares with error bars show mean $\pm$ SEM. (A) In Nutr$-$ medium, where interactions are weak, same-parent and cross-parent pairs show no significant difference in selection correlation, consistent with weak community-level selection. (B) In Nutr+ medium, where interactions are strong, same-parent species exhibit significantly higher selection correlation than cross-parent pairs ($p < 0.001$), indicating correlated species selection that underlies community-level selection.}
\label{fig:suppl_s21}
\end{figure}

% Supple Fig S22 - Linear regression coefficients
\begin{figure}[H]
\centering
\fbox{\parbox{0.9\textwidth}{\centering\vspace{2cm}\textbf{[MISSING] Linear regression coefficients}\vspace{2cm}}}
\caption{Linear coefficients of the fitted regression model showing that a small set of five species accounted for most of the predictive signal in Nutr+ media.}
\label{fig:suppl_s22}
\end{figure}

% Supple Fig S23 - Species-pH correlation
\begin{figure}[H]
\centering
\fbox{\parbox{0.9\textwidth}{\centering\vspace{2cm}\textbf{[MISSING] Species-pH correlation}\vspace{2cm}}}
\caption{Change in monoculture: the relative abundance of these species are highly correlated with the final pH of the parental community.}
\label{fig:suppl_s23}
\end{figure}

% Supple Fig S24 - Hierarchical interactions
\begin{figure}[H]
\centering
\fbox{\parbox{0.9\textwidth}{\centering\vspace{2cm}\textbf{[MISSING] Hierarchical interactions}\vspace{2cm}}}
\caption{Introducing more hierarchical effective interactions.}
\label{fig:suppl_s24}
\end{figure}

% ============================================
% Assembly Effect Figures (from assembly_effect.tex)
% ============================================

% Supple Fig S25 - Assembly reduces interaction strength
\begin{figure}[H]
\centering
\includegraphics[width=0.6\textwidth]{supplementary_figs/Assembly_effect_scatter_combined.png}
\caption{\textbf{Assembly reduces mean pairwise interaction strength.} Mean interaction strength before assembly (red, inter-community) versus after assembly (blue, intra-community) across different interaction strength parameters $\mu$. Points show individual simulation runs; lines connect means. Post-assembly communities consistently show reduced interaction strength, indicating that assembly filters out strongly competing species.}
\label{fig:suppl_s25}
\end{figure}

% Supple Fig S26 - Assembly effect simulation
\begin{figure}[H]
\centering
% TODO: Convert Fig_assembly_effect_simulation.svg to PDF
\fbox{\parbox{0.9\textwidth}{\centering\vspace{2cm}\textbf{[SVG] Fig\_assembly\_effect\_simulation.svg}\vspace{2cm}}}
\caption{\textbf{Assembly history promotes CLS in simulations.} Pie charts comparing outcome distributions between coalescence (top row) and direct assembly (bottom row) across five interaction strengths. Coalescence consistently produces higher CLS (red) and lower Restructuring (purple) compared to direct assembly, demonstrating that assembly history creates correlated species retention.}
\label{fig:suppl_s26}
\end{figure}

% Supple Fig S27 - Assembly effect experimental
\begin{figure}[H]
\centering
% TODO: Convert Fig_assembly_effect_experimental.svg to PDF
\fbox{\parbox{0.9\textwidth}{\centering\vspace{2cm}\textbf{[SVG] Fig\_assembly\_effect\_experimental.svg}\vspace{2cm}}}
\caption{\textbf{Assembly history promotes CLS in experiments.} Pie charts comparing outcome distributions between coalescence (top row) and direct assembly (bottom row) across three nutrient conditions. In Base and Nutr$+$ conditions, coalescence shows elevated CLS rates compared to direct assembly, consistent with simulation predictions.}
\label{fig:suppl_s27}
\end{figure}

% ============================================
% Natural Communities Figure (from natural_communities.tex)
% ============================================

% Supple Fig S28 - Natural communities coalescence
\begin{figure}[H]
\centering
\includegraphics[width=\textwidth]{supplementary_figs/Fig_natural_communities_coalescence.pdf}
\caption{\textbf{Coalescence outcomes in natural sample-derived communities.}
\textbf{(A)} Experimental workflow: natural environmental samples are cultured through serial growth-dilution to establish stable communities, then mixed pairwise for coalescence experiments.
\textbf{(B)} Phase diagrams showing coalescence outcomes in parental similarity space across three nutrient conditions. Each point represents a coalescence event; dashed lines indicate equal similarity to both parents. Histograms below show the distribution of Parental Dominance Index (PDI).
\textbf{(C)} Stacked bar chart summarizing coalescence outcome fractions (CLS: community-level selection, M: Mixing, R: Restructuring) across media conditions. Natural communities exhibit the same CLS-dominated patterns observed in synthetic communities, demonstrating the generality of one-sided selection dynamics.}
\label{fig:suppl_s28}
\end{figure}

% ============================================
% Pairwise Selection Correlation vs Interaction Strength
% ============================================

% Supple Fig S29 - Pairwise selection correlation vs interaction strength
\begin{figure}[H]
\centering
\includegraphics[width=0.7\textwidth]{supplementary_figs/correlation_vs_interaction_strength.pdf}
\caption{\textbf{Pairwise selection correlation increases with interaction strength.} Mean pairwise selection correlation between species pairs across simulated coalescence events as a function of mean interaction strength $\mu$. Species pairs from the same parent community (red circles) show positive correlation that increases with $\mu$, indicating correlated survival. Cross-parent species pairs (blue squares) show negative correlation that becomes more negative with $\mu$, indicating competitive exclusion. The grey horizontal line shows the random selection baseline. Error bars represent standard error across $n = 1200$ coalescence events per interaction strength. This pattern demonstrates that stronger interspecies interactions amplify the assembly effect, leading to increased community-level selection.}
\label{fig:suppl_s29}
\end{figure}

% ============================================
% Monoculture Characterization
% ============================================

% Supple Fig S30 - Monoculture OD and growth rate characterization
\begin{figure}[H]
\centering
\includegraphics[width=0.9\textwidth]{supplementary_figs/monoculture_od_growth_histograms.pdf}
\caption{\textbf{Phenotypic diversity of bacterial isolates in monoculture (Base medium).} (A) Distribution of optical density (OD$_{600}$) reached by 55 bacterial ASVs during monoculture growth in Base medium (mean = 0.45, std = 0.12). (B) Distribution of growth rates (h$^{-1}$) across 60 ASVs (mean = 11.1~h$^{-1}$, std = 3.3~h$^{-1}$). The isolates exhibit broad variation in both growth yield and growth rate, reflecting phenotypic diversity that may contribute to varied coalescence outcomes in community mixing experiments.}
\label{fig:suppl_s30}
\end{figure}

% Supple Fig S31 - Overlap fraction histogram
\begin{figure}[H]
\centering
\includegraphics[width=0.95\textwidth]{supplementary_figs/overlap_fraction_histogram.pdf}
\caption{\textbf{Fraction of coalesced ASVs present in both parental communities.} Distribution of overlap fraction across coalescence events in three nutrient conditions. Overlap fraction is defined as the proportion of ASVs in the coalesced community that were present in both parental communities before mixing. Nutr$-$: mean = 0.10 $\pm$ 0.10 (n=94); Base: mean = 0.22 $\pm$ 0.26 (n=94); Nutr+: mean = 0.17 $\pm$ 0.25 (n=94). The relatively low overlap fractions across all conditions indicate that most surviving ASVs in coalesced communities originated from only one of the two parental communities, consistent with the prevalence of one-sided community-level selection outcomes.}
\label{fig:suppl_s31}
\end{figure}
